% ============================================================
% Section 3.2: Subsystem Architecture Breakdown
% ============================================================
\section{Subsystem Architecture Breakdown}
\label{sec:subsystem_architecture_breakdown}
\label{sec:system_operations_description}

This section details the functional architecture of the four subsystems comprising the integrated EOG-guided assistive platform, detailing the signal transformation pipeline from biopotential sensing to actuation and virtual communication.

% ----------------------------------------------------------
% Subsystem 1: Biopotential Acquisition & AFE
% ----------------------------------------------------------
\subsection{Subsystem 1: Biopotential Acquisition and Analog Front-End (AFE)}
\label{subsec:subsystem1_afe}
\label{subsec:physiological_signal_generation}
\label{subsec:differential_amplification}
\label{subsec:analog_bandpass_filtering}
\label{subsec:dc_level_shifting}

\begin{itemize}
    \item \textbf{Ag/AgCl Surface Electrodes:} Detects the microvolt-level potential differences generated by the standing corneo-retinal dipole during voluntary eye movements and blinks. Electrodes are placed in a vertical periorbital configuration with a forehead reference electrode.
    \item \textbf{Instrumentation Amplifier Stage (AD620):} Provides primary differential amplification for weak raw EOG potentials ($\approx 100\,\mu\text{V} - 600\,\mu\text{V}$). The stage provides a Common-Mode Rejection Ratio ($\text{CMRR} > 100\text{ dB}$) and an input impedance of $10^{10}\,\Omega$ to reject power-line noise and prevent skin-electrode loading effects.
    \item \textbf{Active Filtering Cascade (TL072 Op-Amps):} Consists of a second-order active Sallen-Key High-Pass Filter ($f_c = 0.8\text{ Hz}$) to remove half-cell potential DC offsets and motion baseline drift, followed by a second-order Low-Pass Filter ($f_c = 30\text{ Hz}$) to eliminate high-frequency electromyographic (EMG) muscle artifacts.
    \item \textbf{50 Hz Power-Line Notch Filter:} A dedicated active notch filter sharply attenuates $50\text{ Hz}$ ambient electromagnetic hum coupled from surrounding power infrastructure.
    \item \textbf{DC Level Shifter:} Adds a stable $+2.5\text{ V}$ DC offset to the amplified bipolar signal ($\pm V$), converting it into a purely positive unipolar waveform ($0\text{ V} - 5\text{ V}$) fully compatible with the analog input limits of the microcontroller.
\end{itemize}

% ----------------------------------------------------------
% Subsystem 2: Digitization & Central DSP Pipeline
% ----------------------------------------------------------
\subsection{Subsystem 2: Digitization and Central Signal Processing Pipeline}
\label{subsec:subsystem2_dsp}
\label{subsec:sampling_quantization}
\label{subsec:serial_telemetry}
\label{subsec:serial_data_streaming}
\label{subsec:digital_signal_filtering}
\label{subsec:feature_extraction_classification}

\begin{itemize}
    \item \textbf{Analog-to-Digital Conversion (ADC):} The conditioned analog waveform is sampled by a 10-bit embedded ADC at a continuous sampling frequency ($F_s = 250\text{ Hz}$), quantizing the $0-5\text{ V}$ input into digital values ($0-1023$).
    \item \textbf{Serial Data Streaming:} Digitized biopotential samples are packaged and transmitted to the host environment via UART serial communication at a baud rate of $115200\text{ bps}$.
    \item \textbf{MATLAB DSP Engine:} Applies dynamic moving-average envelope filtering and secondary digital Butterworth notch filtering to yield a clean, smoothed waveform.
    \item \textbf{Feature Extraction \& Peak Detection:} Computes instantaneous signal amplitude against defined threshold levels ($V_{th}$) and measures the pulse time duration ($T_{width}$) to reliably discriminate intentional control blinks from involuntary spontaneous blinks.
\end{itemize}

% ----------------------------------------------------------
% Subsystem 3: System Mode Manager & Output Interfaces
% ----------------------------------------------------------
\subsection{Subsystem 3: System Mode Manager and Output Interfaces}
\label{subsec:subsystem3_mode_manager}
\label{subsec:fsm_decision_engine}
\label{subsec:system_mode_management}
\label{subsec:virtual_keyboard_gui}

\begin{itemize}
    \item \textbf{Finite State Machine (FSM) Mode Manager:} Serves as the central decision engine. It interprets blink counts and temporal sequences. Upon classifying a four-blink trigger sequence, the FSM toggles the active operational context between communication and mobility modes.
    \item \textbf{Virtual Optical Speller (Mode A):} When active, classified blink commands drive a circular 6-sector auto-scanning interface, enabling hands-free Arabic text entry, editing, and display management.
    \item \textbf{Wireless Telemetry (Mode B):} When Wheelchair Mode is active, classified directional commands are converted into control bytes and passed to an HC-05 Bluetooth master module for wireless transmission.
\end{itemize}

% ----------------------------------------------------------
% Subsystem 4: Autonomous Robotic Wheelchair & Safety Override
% ----------------------------------------------------------
\subsection{Subsystem 4: Autonomous Robotic Wheelchair and Safety Override}
\label{subsec:subsystem4_wheelchair_safety}
\label{subsec:smart_wheelchair_control}

\begin{itemize}
    \item \textbf{Embedded Control Receiver:} An HC-05 slave receiver mounted on the prototype wheelchair receives the wireless control frames and forwards them to the navigation microcontroller.
    \item \textbf{Ultrasonic Safety Array (HC-SR04):} Continuously polls front ($40\text{ cm}$ threshold) and rear ($30\text{ cm}$ threshold) distances using Time-of-Flight (ToF) acoustic measurements.
    \item \textbf{Autonomous Safety Override:} If an environmental obstacle breaches the safety perimeter, an immediate hardware interrupt suspends active user EOG drive commands, arresting motor power and triggering an automated escape rotation.
    \item \textbf{H-Bridge Power Stage \& Actuation (L298N \& DC Motors):} Drives four geared DC motors configured in a differential steering layout. The driver applies Pulse-Width Modulation (PWM) signals to ensure smooth acceleration, linear speed regulation, and zero-radius turns.
\end{itemize}
